\section{Contributions}
\label{section:contribution}

\subsection{Core Proof-Theoretic Results}

Listed below are the proof-theoretic results that we have mechanized for our
representative subset of SAX, covering identity ($\texttt{id}$), implication
(${\supset}R$, ${\supset}L^0$), tensor (${\otimes}R^0$, ${\otimes}L$),
and snips ($\texttt{snip}^1$, $\texttt{snip}^2$). 

\begin{theorem}[Translations between $G_3$ and SAX]
  $\Gamma \vdash A$ in $G_3$ iff $\Gamma \vdash A$ in SAX.
\end{theorem}

In the following theorems, we say a derivation is cut-free if it contains
no uses of cut, but may contain snips.

\begin{theorem}[Admissibility of Cut]
  If there are cut-free derivations of $\Gamma \vdash x : A$ and
  $\Gamma, x : A \vdash z : C$, then there exists a cut-free derivation
  of $\Gamma \vdash z : C$.
\end{theorem}

\begin{theorem}[Cut Elimination in SAX]
  If there is a derivation of $\Gamma \vdash x : A$, then there is a
  cut-free derivation of the same sequent.
\end{theorem}

\begin{theorem}[Subformula Property]
  All formulas in a cut-free derivation of $\Gamma \vdash x : A$ are
  subformulas of formulas in $\Gamma$ or $A$.
\end{theorem}


\subsection{Encoding SAX in Beluga}

\subsection{Eligibility and Irrelevance}

\begin{lstlisting}
  LF lconcS : o → type =
    ...
    | l⊃RS : (lhypS A → lconcS B) → lconcS (A ⊃ B)
    | l⊃L0S : lhypS A → lhypS (A ⊃ B) → lconcS B
    ...    
  and
  LF elig_c : lconcS A → type =
    | elig_c/l⊃L0S : {x : lhypS A} {y : lhypS (A ⊃ B)}
                     elig_c (l⊃L0S x y)
  and
  LF elig_h : (lhypS A → lconcS B) → type =
    | elig_h/l⊃L0S : {y : lhypS (A ⊃ B)}
                     elig_h (\x. l⊃L0S x y)
  ;
\end{lstlisting}

\subsection{Translation Between G3 and SAX}

\begin{lstlisting}
  rec seq2sax : (A : [⊢ o]) rel-ctx-S [Γ] [Γ'] → [Γ ⊢ conc A[]] → [Γ' ⊢ concS A[]] =
  / total 2 /
  fn R D ⇒
  case D of
    ...
    | [Γ ⊢ ⊃L c1 (\h. c2) H] ⇒
      let [Γ' ⊢ HS] = hyp2hypS R [Γ ⊢ H] in
      let [Γ' ⊢ c1S] = seq2sax R [Γ ⊢ c1] in
      let [Γ', x : hypS B[] ⊢ c2S] = seq2sax (cons-S R) [Γ, x : hyp _ ⊢ c2] in
      [Γ' ⊢ cutS (cutS c1S (\h. ⊃L0S h HS[..])) (\h. c2S)]
    ...
  ;
\end{lstlisting}

\begin{lstlisting}
  rec sax2seq : (A : [⊢ o]) rel-S-ctx [Γ'] [Γ] → [Γ' ⊢ concS A[]] → [Γ ⊢ conc A[]] =
  / total 2 /
  fn R D ⇒
  case D of
    ...
    | [Γ' ⊢ ⊃L0S HS1 HS2] ⇒
      let [Γ ⊢ H1] = hypS2hyp R [Γ' ⊢ HS1] in
      let [Γ ⊢ H2] = hypS2hyp R [Γ' ⊢ HS2] in
      [Γ ⊢ ⊃L (initial H1) (\h. initial h) H2]
    ...
  ;
\end{lstlisting}

\subsection{Admissibility of Cut and Cut Elimination}

\subsection{Subformula Property}